\documentclass[11pt]{article}
\usepackage{fontspec}
\setmainfont{Times New Roman}
\setsansfont{Arial}
\setmonofont{Consolas}
\usepackage{amsmath,amssymb,mathtools}
\usepackage{geometry}
\usepackage{microtype}
\geometry{a4paper,margin=2.28cm}
\setlength{\parindent}{1.15em}
\setlength{\parskip}{0pt}
\makeatletter
\renewcommand\footnotesize{\@setfontsize\footnotesize{7.5}{8.5}\selectfont}
\makeatother
\providecommand{\mG}{\mathfrak{G}}
\providecommand{\mH}{\mathfrak{H}}
\providecommand{\mA}{\mathfrak{A}}
\providecommand{\mB}{\mathfrak{B}}
\providecommand{\mM}{\mathfrak{M}}
\providecommand{\mo}{\mathfrak{o}}
\providecommand{\OO}{\Omega}
\providecommand{\Th}{\Theta}
\emergencystretch=120em
\allowdisplaybreaks
\sloppy
\hbadness=10000
\vbadness=10000
\hfuzz=2pt
\vfuzz=10pt
\raggedbottom
\begin{document}
% Unit: Paper34 section01 v001. Direct interslavic translation from audited German source lines 49--102, segments P34-S0015--P34-S0027. Source scan pages 5--7 are retained as witness, and Zenodo record 20822156 was checked at 2026-06-24T05:56Z.
\setcounter{footnote}{5}

\section*{I. poglavje. Grupno-teoretične osnovy}

\subsection*{§ 1. Grupy s operatorami}

Nehaj $\mG$ bude grupa (konečna ili beskonečna), s elementami $a,b,\ldots$. Pod oblastju operatorov $\Omega$ za grupu $\mG$ razuměje se množina novih simbolov $H,\Theta,\ldots$, taka že vsakomu $a$ iz $\mG$ i vsakomu $\Theta$ iz $\Omega$ jednoznačno jest opreděljen element $\Theta a$ iz $\mG$, i že zadovoljen jest distributivny zakon:
\[
        \Theta(ab)=\Theta a\cdot\Theta b.
\]
Soglasno tomu vsaky operator opredělja homomorfizm grupy v sebe, pri čem grupa obrazuje se na sebe samu ili na podgrupu. Jediničny element prěhodi v jediničny element, obratny element v obratny.

Oblast operatorov nazyvaje se \emph{absolutnoju}, ako različni operatori takože opreděljajut različne homomorfizmy. Iz libovoljnoj oblasti operatorov dobiva se absolutna tako, že sravnjajut se vsi elementi, ktore poraždajut toj samy homomorfizm. Absolutna oblast operatorov jest jedino-jednoznačny obraz podmnožiny množiny vsih homomorfizmov grupy v sebe.

\emph{Dopustima podgrupa} $\mH$ od $\mG$ -- vzhodno k fiksovanoj oblasti operatorov $\Omega$ -- jest taka podgrupa, ktora dopuskaje tu oblast operatorov, t. j. $\Theta a$ leži v $\mH$ za vsako $a$ iz $\mH$ i vsako $\Theta$ iz $\Omega$.

\emph{Priměry.} 1. Nehaj operatori budut vnutrne avtomorfizmy: $\Theta a=c^{-1}ac$. Dopustime podgrupy sut normalne dělitelje.

2. Nehaj operatori budut vsi avtomorfizmy. Dopustime podgrupy sut “karakteristične podgrupy”, ktore pri vsih avtomorfizmah prěhodet v sebe.\footnote{Tut objasnjene pojmy izhode od W. Krull, \emph{Über verallgemeinerte endliche Abelsche Gruppen}, Math. Zeitschr. 23 (1925), s. 161--196, i O. Schmidt, \emph{Über unendliche Gruppen mit endlicher Kette}, Math. Zeitschr. 29 (1928), s. 34--41.}

3. Nehaj $\mG$ bude kolco, t. j. Abelova grupa vzhodno k dodavanju, gde takože opreděljeno jest množenje s vlastnostjami
\[
        r(a+b)=ra+rb,\qquad (a+b)r=ar+br,\qquad ab\cdot c=a\cdot bc.
\]
Vsaky element $r$ Jednočasno opredělja dva operatora: operatori $rx$ i $xr$. Dopustime podgrupy sut “idealy” $\mathfrak a$, imenno: lěvostrannje, ktore dopuskajut operacije $rx$: $r\mathfrak a\subseteq\mathfrak a$; pravostrannje, ktore dopuskajut operacije $xr$: $\mathfrak a r\subseteq\mathfrak a$; dvustrannje, ktore dopuskajut obě operacije. Vsi idealy stavajut se trivialne $(=\{0\}\text{ ili }=\mG)$, ako kolco $\mG$ jest tělo, t. j. ako $\mG-\{0\}$ jest grupa vzhodno k množenju.\footnote{Tedy ni za kolca ni za těla ne predpostavlja se komutativny zakon množenja.}

4. \emph{Moduly vzhodno k kolcu $\mo$.} Nehaj $\mo$ bude kolco, a $\mM$ Abelova grupa, pisana aditivno. Nehaj dano jest množenje $r\cdot a$ elementov iz $\mo$ s elementami iz $\mM$; produkt znovu imaje ležati v $\mM$. Trěba požadovati:
\[
\left.
\begin{aligned}
 r(a+b)&=ra+rb,\\
 rs\cdot a&=r\cdot sa,\\
 (r+s)a&=ra+sa,
\end{aligned}
\right\}\qquad r,s\in\mo,\quad a,b\in\mM.
\]
Togda $\mM$ nazyvaje se $\mo$-modulom; točnije, poneže množitelje iz $\mo$ pišut se lěvo, lěvym $\mo$-modulom. Kolco $\mo$ jest Jednočasno oblast operatorov. Ako ona jest absolutna, t. j. ako različni elementi kolca daju takože različne operacije, govori se o \emph{absolutnoj oblasti množiteljev} za modul $\mM$.

Kako vyše, iz libovoljnoj oblasti množiteljev možno prějdti k absolutnoj tako, že sravnjajut se elementi, ktore poraždajut tu samu operaciju. Absolutna oblast množiteljev jest znovu kolco. Dopustime podgrupy modula $\mM$ nazyvajut se \emph{podmoduly}. Ako specialno $\mo=\mM$, togda vrčajemo se k lěvym idealam.\footnote{Tym samym načinom za prave moduly požada se:
\[
(a+b)r=ar+br,\qquad a\cdot rs=ar\cdot s,\qquad a(r+s)=ar+as.
\]}

5. \emph{Dvomoduly.} Ako $\mM$ jest Jednočasno lěvy $\mo$-modul i pravy $\mo'$-modul, i ako nadto za $a$ v $\mo$, $\alpha$ v $\mM$, $a'$ v $\mo'$ vsegda
\[
        a\cdot\alpha a'=a\alpha\cdot a',
\]
togda $\mM$ nazyvaje se dvomodulom.

6. \emph{Kolco avtomorfizmov Abelovej grupy.}\footnote{Sr. A. Châtelet, \emph{Les Groupes Abéliens finis}, Paris, Gauthier-Villars, 1925, p. 99.} Za homomorfizmy v sebe jednej Abelovej grupy (pisanoj aditivno) možno opreděliti dodavanje i množenje formulami
\[
        (H+\Theta)a=Ha+\Theta a,\qquad (H\Theta)a=H(\Theta a).
\]
Tako opreděljene operacije javno znovu predstavjajut homomorfizmy i zato znovu prinadležet tomu sistemu. Sistem tvori kolco, i Abelovu grupu možno po 4. smatrjati kako modul vzhodno k tomu kolcu.

Pod “podgrupami” dalje vsegda razumějut se dopustime podgrupy. Prěsěk $\mA\cap\mB$ dvoh dopustimyh podgrup jest znovu dopustima podgrupa. Takože produkt $\mA\mB$, ako najmanje jedna iz dvoh podgrup $\mA,\mB$ jest normalny dělitelj.
\end{document}
